% !TEX program = xelatex
% 第二章总技术路线图：独立可编译的 TikZ 源文件。
% 变量沿用当前论文：领域配比记为 p；质量/配比跨附件衔接为条件情景，
% 损失到能力的映射需传播误差。本图展示研究流程，不代表各环节均已验证。
\documentclass[UTF8,a4paper,10pt]{ctexart}
\usepackage[margin=9mm]{geometry}
\usepackage{amsmath}
\usepackage{tikz}
\usetikzlibrary{arrows.meta,calc,positioning}
\pagestyle{empty}
\setlength{\parindent}{0pt}

\tikzset{
  flow/.style={-{Stealth[length=2mm,width=1.35mm]},line width=.43pt},
  transfer/.style={-{Stealth[length=2.7mm,width=1.8mm]},line width=1.05pt},
  conditional/.style={-{Stealth[length=2mm,width=1.35mm]},
    densely dashed,line width=.72pt},
  module/.style={draw=black,line width=.38pt,fill=white,align=center,
    inner xsep=1.4mm,inner ysep=.65mm,minimum height=7mm,
    font=\fontsize{7.5}{8.7}\selectfont},
  core/.style={module,double,double distance=.52pt,line width=.44pt},
  result/.style={module,line width=.76pt},
  region/.style={draw=black,densely dashed,line width=.5pt,rounded corners=0pt},
  regiontitle/.style={fill=white,inner xsep=1.2mm,inner ysep=.2mm,
    font=\bfseries\fontsize{9}{10}\selectfont},
  smalllabel/.style={font=\fontsize{7}{8}\selectfont,fill=white,inner sep=.4mm}
}

\begin{document}
\noindent\begin{tikzpicture}[x=1mm,y=1mm]
\useasboundingbox (0,0) rectangle (192,278);

% 目标层
\node[font=\bfseries\fontsize{15}{17}\selectfont] at (96,270) {总技术路线图};
\node[core,minimum width=145mm,minimum height=12mm,
  font=\bfseries\fontsize{10}{11.5}\selectfont] (goal) at (96,255)
  {有限算力下的大语言模型能力提升与资源配置};

% 问题一：质量与领域结构
\draw[region] (3,185) rectangle (189,246);
\node[regiontitle,anchor=west] at (9,246) {问题一\quad 数据质量评价与领域配比分析};
\node[module,minimum width=116mm] (q1in) at (96,235) {附件 A 训练语料\quad 质量指标\quad 配比实验};
\node[module,minimum width=116mm] (q1pre) at (96,225) {数据清洗\quad 指标标准化\quad 样本/扩展数据划分};
\node[module,minimum width=74mm,minimum height=10mm] (q1quality) at (53,211)
  {质量评价：综合评分、冲突分析\\排序检验、扩展数据一致性};
\node[module,minimum width=74mm,minimum height=10mm] (q1mix) at (139,211)
  {领域配比：单域、组合效应\\交互作用与配比特征};
\node[core,minimum width=97mm] (q1loss) at (96,198) {质量—配比—相对训练损失关系};
\node[result,minimum width=97mm] (q1out) at (96,190) {$Q_A$、$\boldsymbol p$ 与相对损失效应\quad｜\quad 检验集验证};
\draw[flow] (goal.south) -- (q1in.north);
\draw[flow] (q1in) -- (q1pre);
\draw[flow] (q1pre.south) -- ++(0,-2) -| (q1quality.north);
\draw[flow] (q1pre.south) -- ++(0,-2) -| (q1mix.north);
\draw[flow] (q1quality.south) -- ++(0,-1) -| (q1loss.north);
\draw[flow] (q1mix.south) -- ++(0,-1) -| (q1loss.north);
\draw[flow] (q1loss) -- (q1out);

% 问题二：跨附件条件衔接与广义标度律
\draw[region] (3,122) rectangle (189,182);
\node[regiontitle,anchor=west] at (9,182) {问题二\quad 多因素广义标度律};
\node[module,minimum width=116mm] (q2in) at (96,172)
  {问题一的 $Q_A,\boldsymbol p$\quad +\quad 附件 B 训练记录};
\node[module,minimum width=116mm] (q2vars) at (96,162)
  {变量构造与口径校准：$N,D,Q,\boldsymbol p,L$};
\node[core,minimum width=106mm,minimum height=9mm] (q2model) at (96,150)
  {广义标度律\quad $\widehat L=f(N,D,Q,\boldsymbol p)$};
\node[module,minimum width=76mm] (q2fit) at (53,138)
  {参数估计\quad 轨迹/跨模型验证};
\node[module,minimum width=76mm] (q2check) at (139,138)
  {残差分析\quad 替代关系与外推边界};
\node[result,minimum width=106mm] (q2out) at (96,128)
  {有适用边界的条件损失预测};
\draw[transfer] (q1out.south) -- (q2in.north);
\draw[flow] (q2in) -- (q2vars);
\draw[flow] (q2vars) -- (q2model);
\draw[flow] (q2model.south) -- ++(0,-1) -| (q2fit.north);
\draw[flow] (q2model.south) -- ++(0,-1) -| (q2check.north);
\draw[flow] (q2fit.south) -- ++(0,-1) -| (q2out.north);
\draw[flow] (q2check.south) -- ++(0,-1) -| (q2out.north);

% 问题三：训练优化主线
\draw[region] (3,19) rectangle (93,119);
\node[regiontitle,anchor=west] at (8,119) {问题三\quad 训练优化主线};
\node[module,minimum width=75mm] (q3in) at (48,108)
  {广义标度律\quad +\quad 算力成本模型};
\node[module,minimum width=25mm] (costtrain) at (18,96) {基础训练};
\node[module,minimum width=25mm] (costquality) at (48,96) {质量提升};
\node[module,minimum width=25mm] (costcontext) at (78,96) {长上下文};
\node[module,minimum width=72mm] (budget) at (48,84) {算力预算约束};
\node[core,minimum width=78mm,minimum height=15mm] (opt) at (48,68)
  {约束优化\quad $\min\widehat L$\\[-.3mm]
   $\scriptstyle C_{\rm train}+C_{\rm quality}+C_{\rm context}\leq C_{\rm budget}$};
\node[module,minimum width=72mm] (solve) at (48,52)
  {优化求解\quad 预算/质量成本敏感性};
\node[result,minimum width=72mm] (q3out) at (48,40)
  {条件配置 $(N^*,D^*,Q^*,\boldsymbol p^*)$、$\widehat L^*$};
\node[module,minimum width=72mm] (critical) at (48,28)
  {成本平衡分析\quad $\to$\quad 临界上下文长度};
\draw[transfer] (48,122) -- (q3in.north);
\draw[flow] (q3in.south) -- (48,102);
\draw[flow] (18,102) -- (costtrain.north);
\draw[flow] (48,102) -- (costquality.north);
\draw[flow] (78,102) -- (costcontext.north);
\draw[flow] (costtrain.south) -- ++(0,-2) -| (budget.north);
\draw[flow] (costquality) -- (budget);
\draw[flow] (costcontext.south) -- ++(0,-2) -| (budget.north);
\draw[flow] (budget) -- (opt);
\draw[flow] (opt) -- (solve);
\draw[flow] (solve) -- (q3out);
\draw[flow] (q3out) -- (critical);

% 问题四：能力演进主线；由问题二与附件 C 共同驱动
\draw[region] (99,19) rectangle (189,119);
\node[regiontitle,anchor=west] at (104,119) {问题四\quad 能力演进主线};
\node[module,minimum width=75mm] (q4in) at (144,108)
  {附件 C 能力数据\quad +\quad 问题二损失/规模};
\node[module,minimum width=75mm] (q4pre) at (144,97)
  {历史记录整理\quad 指标统一};
\node[module,minimum width=35mm] (bridge) at (121,84) {损失—能力映射};
\node[module,minimum width=35mm] (history) at (167,84) {历史能力演进};
\node[module,minimum width=35mm] (bridgeerr) at (121,71) {映射误差分析};
\node[module,minimum width=35mm] (contrib) at (167,71)
  {规模关联\\非规模综合项};
\node[core,minimum width=75mm] (frontier) at (144,58) {能力前沿演进模型};
\node[module,minimum width=75mm] (future) at (144,46) {未来 12 / 24 个月情景预测};
\node[module,minimum width=75mm] (uncertainty) at (144,34) {预测不确定性与适用边界};
\node[result,minimum width=75mm] (q4out) at (144,26) {条件能力前沿及情景范围};
\draw[transfer] (144,122) -- (q4in.north);
% Q4 的历史分析可直接使用 C；Q3 条件 Loss 仅经带误差桥接进入能力解释。
\draw[conditional] (q3out.east) -- (96,40) -- (96,84) -- (bridge.west);
\node[smalllabel,rotate=90] at (96,61) {Q3 条件损失};
\draw[flow] (q4in) -- (q4pre);
\draw[flow] (q4pre.south) -- ++(0,-1) -| (history.north);
\draw[flow] (q4pre.south) -- ++(0,-1) -| (bridge.north);
\draw[flow] (history) -- (contrib);
\draw[flow] (bridge) -- (bridgeerr);
\draw[flow] (contrib.south) -- ++(0,-1) -| (frontier.north);
\draw[flow] (frontier) -- (future);
\draw[flow] (future) -- (uncertainty);
\draw[flow] (bridgeerr.south) -- ++(0,-2) -- (103,64) -- (103,34) -- (uncertainty.west);
\draw[flow] (uncertainty) -- (q4out);

% 两条主线在结论层汇合；Q3 的条件 Loss 经桥接进入 Q4 的能力解释。
\node[result,minimum width=132mm,minimum height=10mm,
  font=\bfseries\fontsize{8.7}{10}\selectfont] (final) at (96,8)
  {有限算力下的资源配置与未来能力发展结论};
\draw[flow] (critical.south) -- (48,17) -- (96,17) -- (final.north);
\draw[flow] (q4out.south) -- (144,17) -- (96,17) -- (final.north);

\end{tikzpicture}
\end{document}
